\documentclass[11pt]{article}
\usepackage[a4paper,margin=27mm]{geometry}
\usepackage[T1]{fontenc}
\usepackage{lmodern,microtype,amsmath,amssymb,amsthm,mathtools,enumitem}
\usepackage[hidelinks,hypertexnames=false]{hyperref}
\newtheorem{theorem}{Theorem}
\newtheorem{lemma}[theorem]{Lemma}
\newtheorem{proposition}[theorem]{Proposition}
\newtheorem{corollary}[theorem]{Corollary}
\theoremstyle{remark}\newtheorem{remark}[theorem]{Remark}
\newcommand{\R}{\mathbb R}
\newcommand{\Dscr}{\mathcal D}
\newcommand{\Escr}{\mathcal E}
\newcommand{\Kscr}{\mathcal K}
\newcommand{\supp}{\operatorname{supp}}
\newcommand{\one}{\mathbf 1}
\title{Edge Dichotomy and Leverage}
\author{}\date{9 October 2026}
\begin{document}\maketitle

\noindent
This note continues \emph{Closed Choice for Independent Rankings} (CC) and
\emph{Closed Choice Beyond Ordinary Integrals} (CBOI). The open question there
is whether (G), (IB) and (S) force a strictly positive finitely additive weak
extension of the comparisons that the obstruction uses. That question is still
open, but the space for a counterexample is now much smaller.
\begin{enumerate}[leftmargin=*,itemsep=2pt]
\item The obstruction needs an extension only for comparisons from coherent
undominated models to regular models (Proposition~\ref{prop:one-sided}).
\item Such an extension exists exactly when no \emph{leverage sequence} exists:
combinations of accepted comparisons whose net effect is a huge improvement on
a set of fixed measure at the cost of a bounded worsening
(Theorem~\ref{thm:leverage}).
\item Under (G) and (S), an additive primary criterion must be either strict
or completely blind on mixed problems (Theorem~\ref{thm:dichotomy}).
Consequently no lexicographic hierarchy of ordinary additive criteria
escapes the obstruction, unless every level ignores mixed problems
(Theorem~\ref{thm:lex}).
\item If two of three worlds are additive, the third cannot loosen the face
where the additive worlds meet: that face is forced to the vacuous trade-off
(Theorem~\ref{thm:one-world}).
\end{enumerate}
An independent line-by-line check found no errors in
Theorems~\ref{thm:dichotomy}, \ref{thm:lex} and~\ref{thm:one-world}; its
minor corrections are incorporated. Theorem~\ref{thm:leverage},
Remark~\ref{rem:weak-star} and Corollary~\ref{cor:reduction} have not had an
independent check.

\section{Setting}

We use the setting of v3, CC and CBOI without change. Worlds are
$1,\ldots,n$. Binary problems are represented by $g\in H=\{g_1>0\}$, and
$\mu$ is Gaussian measure normalized on $H$. Write $H'$ for the mixed problems
and $U$ for the unanimous ones. Guidance ranks at world $i$ are $0,1,2,3$
(good, silent, bad, contradictory). A model is \emph{regular} if it never
contradicts itself and is correct on unanimous problems. Coherent models are
regular. For distinct worlds $i,j$, put
\[
 D_{ij}\coloneqq\{g\in H':g_ig_j<0\},\qquad
 t\coloneqq|g_i|/|g_j|\ \text{on }D_{ij}.
\]
The image of $\mu|_{D_{ij}}$ under $t$ has a positive density on $(0,\infty)$.
For $n\ge3$, $H'=\bigcup_{j\ne i}D_{ij}$, and $D_{ij}\cap D_{ik}$ is non-null
for distinct $j,k\ne i$.

The \emph{one-sided comparisons} at world $i$ are
\[
 \Kscr_i\coloneqq\{(C,X):C\text{ coherent and undominated},\ X\text{ regular},\
 C\preceq_iX\}.
\]
We use the following consequences of (LA) and (S):
\begin{enumerate}[label=(KKM\alph*),leftmargin=*]
\item every model is weakly dominated at all worlds by a selected $C_q$;
\item every undominated model is tied at every world with a selected $C_p$
and is competitive on $\supp p$.
\end{enumerate}
By CC's Lemma~7 (face support), a selected model on an edge $[e_i,e_j]$ has
credal set $\{se_i+(1-s)e_j:s\in[a,b]\}$. Put
$\alpha\coloneqq(1-a)/a\ge\beta\coloneqq(1-b)/b$. Orient each problem of
$D_{ij}$ so that $g_i>0>g_j$. Then the selected model:
\begin{itemize}[leftmargin=*,itemsep=1pt]
\item accepts $g$ (good at $i$, bad at $j$) where $t>\alpha$;
\item is silent where $\beta<t<\alpha$;
\item accepts $-g$ (bad at $i$, good at $j$) where $t<\beta$.
\end{itemize}
Off $D_{ij}$ it is correct at both $i$ and $j$. The sets $\{t=\alpha\}$ and
$\{t=\beta\}$ are null.

\section{One-sided extensions and leverage}

\begin{proposition}\label{prop:one-sided}
In CC's Theorem~1 and CBOI's Theorem~2 it suffices that the weak extension
inequality $L_i(C)\le L_i(X)$ hold for $(C,X)\in\Kscr_i$.
\end{proposition}
\begin{proof}
Both proofs compare only two kinds of pair. The first is a selected coherent
model with a regular model: $C_q\preceq_iB$ with $B$ the pointwise Bayes
model, or $C_m\preceq_iX_m$. The second is the target, assumed undominated and
coherent, with a regular model: $\Dscr_M\preceq_j\Escr^\circ$, or
$T\preceq_iX_m$. Selected models are coherent and undominated.
\end{proof}

Fix a world $i$. Let $\Omega'=H'\times\{1,2\}$ carry the probability
$P\coloneqq\frac12(\mu|_{H'}/\mu(H'))\otimes(\text{counting})$. Encode a
regular model $\Dscr$ by
\[
 u^\Dscr\coloneqq\bigl(\one\{r_i(\cdot,\Dscr)\ge1\},\ \one\{r_i(\cdot,\Dscr)\ge2\}\bigr)
 \in L^\infty(\Omega'),
\]
and let $K_i$ be the convex cone generated by $u^X-u^C$, $(C,X)\in\Kscr_i$.

\begin{theorem}\label{thm:leverage}
The following are equivalent.
\begin{enumerate}[label=(\alph*)]
\item There are finite nonnegative finitely additive $\nu_1,\nu_2$ on $H'$, with
$\nu_k(A)=0$ if and only if $\mu(A)=0$, such that
\[
 \nu_1\{r_i(\cdot,X)\ge1\}-\nu_1\{r_i(\cdot,C)\ge1\}
 +\nu_2\{r_i(\cdot,X)\ge2\}-\nu_2\{r_i(\cdot,C)\ge2\}\ \ge\ 0
\]
for every $(C,X)\in\Kscr_i$.
\item There is no \emph{leverage sequence}: no $\alpha>0$ such that for every
$n$ some $\kappa\in K_i$ satisfies $\kappa\le1$ a.e.\ and
$P\{\kappa\le-n\}\ge\alpha$.
\end{enumerate}
In case (a), $\nu_{i1}\coloneqq\nu_1+\mu|_U$, $\nu_{i2}\coloneqq\nu_2+\mu|_U$
and $\nu_{i3}\coloneqq\mu$ define a loss of CBOI's form that satisfies the
inequality of Proposition~\ref{prop:one-sided} at world $i$.
\end{theorem}
\begin{proof}
\emph{(a) implies (b).} Combine $\nu_1,\nu_2$ into one charge $\hat\nu$ on
$\Omega'$. The pairing in (a) says $\hat\nu(\kappa)\ge0$ for every
$\kappa\in K_i$.

By the Yosida--Hewitt decomposition, $\hat\nu=\hat\nu_c+\pi$ with
$\hat\nu_c=f\,P$ countably additive and $\pi$ purely finitely additive.
As in CBOI's Lemma~3, the concentration of $\pi$ on sets of arbitrarily small
measure and the hypothesis $\hat\nu(A)>0$ for non-null $A$ force $f>0$ a.e.
Choose $\delta>0$ with $P\{f<\delta\}<\alpha/2$. Then $\hat\nu(S)\ge\delta\alpha/2$
whenever $P(S)\ge\alpha$.

Given a leverage sequence, let $S_n=\{\kappa_n\le-n\}$. Then
\[
 \hat\nu(\kappa_n)\le\hat\nu(\Omega')-n\,\hat\nu(S_n)
 \le\hat\nu(\Omega')-n\delta\alpha/2,
\]
which is negative for large $n$. This is a contradiction.

\emph{(b) implies (a).} Let $\mathcal C\coloneqq-\overline{K_i+L^\infty_+}$,
with closure in the norm of $L^\infty(\Omega')$. It is a norm-closed convex
cone containing $-L^\infty_+$.

We claim $\mathcal C\cap L^\infty_+=\{0\}$. Suppose $0\ne y\in\mathcal C$ with
$y\ge0$. Then $y\ge\delta\one_S$ for some $\delta>0$ and non-null $S$. For each
$n$, choose $\kappa\in K_i$ and $w\ge0$ with
$\|\kappa+w+y\|_\infty<\delta/(n+1)$. Then
\[
 \kappa\le-\delta\one_S+\frac{\delta}{n+1}\quad\text{a.e.},
 \qquad\text{so}\qquad
 \lambda\kappa\le-(n+1)\one_S+1\quad\text{with }\lambda=\frac{n+1}{\delta}.
\]
Thus $\lambda\kappa\in K_i$ satisfies $\lambda\kappa\le1$ and
$\lambda\kappa\le-n$ on $S$. This is a leverage sequence with $\alpha=P(S)$,
contrary to (b).

Rokhlin's Kreps--Yan theorem for $L^\infty$ \cite{Rokhlin} now gives a strictly
positive $\hat\nu\in(L^\infty)^*$, meaning $\hat\nu(A)>0$ whenever $P(A)>0$,
with $\hat\nu\le0$ on $\mathcal C$. Hence $\hat\nu\ge0$ on $K_i$. Its
restrictions to the two copies of $H'$ give $\nu_1,\nu_2$. Elements of
$(L^\infty)^*$ vanish on null sets, so (a) holds.

\emph{The completed loss.} For regular $C,X$, the unanimous terms and the
rank-$3$ term agree. So $L_i(X)-L_i(C)$ equals the expression in (a).
Strict positivity of $\nu_{ik}$ on all of $H$ holds by construction.
\end{proof}

Proposition~\ref{prop:one-sided}, Theorem~\ref{thm:leverage} and CBOI's
Theorem~2 together give: \emph{if $n\ge3$, (G), (LA) and (S) hold and the
vacuous model is undominated, then some $\Kscr_i$ has a leverage sequence.}

The simplest leverage sequence looks like this. Let $C$ be coherent and
undominated, and let $A\subseteq H'$ be non-null. Suppose that for pairwise
disjoint sets $B_1,B_2,\ldots\subseteq H'\setminus A$ there are regular $X_m$
with $C\preceq_iX_m$, where $X_m$ lowers the rank of $C$ at $i$ by one on $A$,
agrees with $C$ off $A\cup B_m$, and raises it on $B_m$. Fix the rank level
$k\in\{1,2\}$ that the improvement crosses on a non-null part $A'$ of $A$.
Then $\kappa_N\coloneqq\sum_{m\le N}(u^{X_m}-u^C)$ satisfies $\kappa_N\le1$
and $\kappa_N=-N$ on $A'\times\{k\}$. Many small worsenings each outweigh the
same fixed improvement.

\begin{remark}\label{rem:weak-star}
The weak$^*$ closure, for $\sigma(L^\infty,L^1)$, plays the same role for
\emph{ordinary} extensions, meaning (a) with $\nu_k=f_k\mu$ and $f_k>0$ a.e.
Such an extension exists if and only if
$-\one_S\notin\overline{K_i+L^\infty_+}^{\,w^*}$ for every non-null
$S\subseteq\Omega'$.

\emph{Proof.} Necessity: $\kappa\mapsto\int\kappa f\,dP$ is
weak$^*$-continuous and nonnegative on $K_i+L^\infty_+$, so also on its
weak$^*$ closure, and it is negative at $-\one_S$.

Sufficiency: for each non-null $S$, separate $-\one_S$ from the weak$^*$-closed
cone. The dual of $(L^\infty,\sigma(L^\infty,L^1))$ is $L^1$, so this gives
$h_S\in L^1$ that is nonnegative on the cone, hence $h_S\ge0$, with
$\int h_S\one_S\,dP>0$. The set of $h\in L^1_+$ nonnegative on $K_i$ is
closed under countable sums $\sum_m2^{-m}h_m/(1+\|h_m\|_1)$, by dominated
convergence. Take such a sum $h^*$ with $P\{h^*>0\}$ maximal. If
$S=\{h^*=0\}$ were non-null, then $h^*+h_S$ would have larger support. So
$h^*>0$ a.e.
\end{remark}

CBOI's Proposition~5 lies between the two closures. Its own loss has the
form of (a), so it has no leverage sequence. It also has no ordinary
extension, even on $\Kscr_i$. To see this, choose $p\ne e_i$ on an edge
$[e_i,e_j]$ so near $e_i$ that $0<L_i(C_p)<a$, in CBOI's notation. Then
$C_p\prec_iB_m$ for every $m$. An ordinary strictly guidance-ordered
$\Lambda_i$ would give
$\Lambda_i(C_p)\le\Lambda_i(B_m)\to\Lambda_i(D_i)<\Lambda_i(C_p)$.
So in CBOI's example some $-\one_S$ lies in the weak$^*$ closure of
$K_i+L^\infty_+$, but no $-\one_S$ lies in its norm closure.

\section{Regular primaries and the edge dichotomy}

A \emph{regular primary} at world $i$ is a real function $\Phi$ on regular
models such that, for regular $\Dscr,\Escr$,
\[
 \Phi(\Dscr)<\Phi(\Escr)\ \Longrightarrow\ \Dscr\prec_i\Escr.
\]
It weakly respects $\preceq_i$ on regular models. Restricting a primary in
the sense of CC or CBOI to regular models gives a regular primary.

A regular primary is \emph{additive} if, for regular $X$,
\[
 \Phi(X)=\int_H\varphi\bigl(g,r_i(g,X)\bigr)\,d\mu(g),
\]
where $\varphi(\cdot,0)\le\varphi(\cdot,1)\le\varphi(\cdot,2)$ are measurable
and bounded in absolute value by an integrable $h$. On $H'$ put
\[
 c\coloneqq\varphi(\cdot,1)-\varphi(\cdot,0)\ge0,\qquad
 d\coloneqq\varphi(\cdot,2)-\varphi(\cdot,1)\ge0.
\]
Then $|\Phi(X)-\Phi(Y)|\le2\int_{\{a_X\ne a_Y\}}h\,d\mu$, so $\Phi$ is
continuous in measure on regular models.

The next lemma restates CC's Lemma~10(a) and CBOI's Lemma~6 for regular
primaries. The proofs are the same, because they compare selected models
only with regular models. We give them in full.

\begin{lemma}\label{lem:regular-primary}
Assume (G) and (S). Let $\Phi$ be a regular primary at $i$, and put
$f(p)\coloneqq\Phi(C_p)$.
\begin{enumerate}[label=(\alph*)]
\item If a regular $\Escr$ satisfies $\Escr\prec_iC_p$ and $\Phi(\Escr)=f(p)$,
then $p$ is a local minimum of $f$ on $\Delta$.
\item Let $p\in[e_i,e_j]$. If a regular $F$ satisfies $F\prec_jC_p$ and
$\Phi(F)=f(p)$, then $p$ is a local maximum of $f$ on $[e_i,e_j]$.
\item If $\Phi$ is continuous in measure on regular models, then $f$ is
continuous on every edge $[e_i,e_j]$.
\end{enumerate}
\end{lemma}
\begin{proof}
(a) Since $C_p\not\preceq_i\Escr$, the point $p$ lies outside the closed set
$A_i(\Escr)$. So some neighbourhood $N$ of $p$ misses it. For $q\in N$, if
$f(q)<\Phi(\Escr)$ then $C_q\prec_i\Escr$, which is false. So
$f(q)\ge\Phi(\Escr)=f(p)$.

(b) Similarly, a neighbourhood $N$ of $p$ in the edge misses $A_j(F)$. Let
$q\in N$. Then $C_q\not\preceq_jF$. Since $C_q$ is competitive on
$\supp q\subseteq\{i,j\}$, we get $C_q\preceq_iF$. So $f(q)\le\Phi(F)=f(p)$.

(c) By CC's Lemma~7 the selections on the edge are the models described in
Section~1. In particular $C_{e_i}=D_i$ and $C_{e_j}=D_j$. By (G),
$D_i\preceq_iX$ for every $X$, so $\Phi(D_i)\le\Phi(X)$ for regular $X$.

\emph{Lower semicontinuity.} Suppose $p_m\to p$ on the edge and
$f(p_m)\le f(p)-\epsilon$ with $\epsilon>0$. Put
$G_s\coloneqq\{g\in H:N(g_2)<s\}$, where $N$ is the standard normal
distribution function. Then $\mu(G_s)=s$. Let $X_s$ follow $C_p$ on $G_s$
and $D_i$ elsewhere. Each $X_s$ is regular, and $s\mapsto\Phi(X_s)$ is
continuous, from $\Phi(D_i)\le f(p)-\epsilon$ to $f(p)$. Choose $s$ with
$\Phi(X_s)=f(p)-\epsilon/2$ and put $\Escr\coloneqq X_s$. Then
$C_{p_m}\prec_i\Escr$ for every $m$, so $p_m\in A_i(\Escr)$, and by
closedness $p\in A_i(\Escr)$. But $\Escr\prec_iC_p$. This is a contradiction.

\emph{Upper semicontinuity.} Suppose $p_m\to p$ and
$f(p_m)\ge f(p)+\epsilon$. Every edge model has rank at most $2$ at $i$ on
$D_{ij}$ and rank $0$ elsewhere, while $D_j$ has rank $2$ on $D_{ij}$. So by
(G), $C_q\preceq_iD_j$ and $f(q)\le\Phi(D_j)$ for every $q$ on the edge.
Hence $C_p\ne D_j$; otherwise $f(p)$ would be the maximum.

So $C_p$ is not good at $j$ on a non-null $W\subseteq D_{ij}$. Choose a
non-null $Z\subseteq W$ so small that the model $\Escr$, which agrees with $C_p$
except that it gives the $j$-good advice on $Z$, has
$\Phi(\Escr)<f(p)+\epsilon$. Then $\Escr$ is regular and $\Escr\prec_jC_p$ by
strict (G). Thus $p\notin A_j(\Escr)$, and for large $m$,
$C_{p_m}\not\preceq_j\Escr$. Competitiveness gives $C_{p_m}\preceq_i\Escr$, so
$f(p_m)\le\Phi(\Escr)<f(p)+\epsilon$. This is a contradiction.
\end{proof}

\begin{theorem}\label{thm:dichotomy}
Assume (G) and (S). Let $\Phi$ be an additive regular primary at world $i$,
and let $j\ne i$. Then exactly one of the following holds:
\begin{itemize}[leftmargin=*,itemsep=1pt]
\item \emph{strict:} $c>0$ and $d>0$ a.e.\ on $D_{ij}$;
\item \emph{blind:} $c=d=0$ a.e.\ on $D_{ij}$.
\end{itemize}
If $n\ge3$, the same alternative holds for every $j\ne i$, and hence a.e.\ on
$H'$.
\end{theorem}
\begin{proof}
\emph{Step 1: the values on the edge.} Let $p\in[e_i,e_j]$ have thresholds
$\alpha\ge\beta$. By Section~1, $C_p$ has rank $0$ at $i$ off $D_{ij}$. On
$D_{ij}$ its rank is $0$ where $t>\alpha$, $1$ where $\beta<t<\alpha$, and $2$
where $t<\beta$. With $m_0\coloneqq\Phi(D_i)$,
\[
 f(p)=m_0+\int_{D_{ij}\cap\{t<\alpha\}}c\,d\mu
 +\int_{D_{ij}\cap\{t<\beta\}}d\,d\mu.
\]
For $x\in[0,\infty]$ put $G(x)\coloneqq\int_{D_{ij}\cap\{t<x\}}(c+d)\,d\mu$.
Since $c,d\ge0$,
\begin{equation}
 m_0+G(\beta)\le f(p)\le m_0+G(\alpha).
 \label{eq:sandwich}
\end{equation}
Also $f(e_i)=m_0$ and $f(e_j)=m_0+G(\infty)$. By
Lemma~\ref{lem:regular-primary}(c) and the intermediate value theorem, $f$
maps the edge onto $[m_0,m_0+G(\infty)]$.

\emph{Step 2: two mechanisms.} For an interval $I\subseteq(0,\infty)$ put
\[
 \kappa^c(I)\coloneqq\mu\bigl(D_{ij}\cap\{c=0\}\cap\{t\in I\}\bigr),\qquad
 \kappa^d(I)\coloneqq\mu\bigl(D_{ij}\cap\{d=0\}\cap\{t\in I\}\bigr).
\]
\begin{itemize}[leftmargin=*,itemsep=2pt]
\item \emph{Minimum.} Suppose $\kappa^c(\beta,\alpha)>0$. Let $Z$ be the set
it measures. There $C_p$ is silent. Let $\Escr$ give the $i$-good advice on
$Z$ and agree with $C_p$ elsewhere. Then $\Escr$ is regular,
$\Escr\prec_iC_p$ by strict (G), and $\Phi(\Escr)=f(p)-\int_Zc=f(p)$.

Suppose instead $\kappa^d(0,\beta)>0$, with $Z$ the set it measures. There
$C_p$ is bad at $i$. Let $\Escr$ be silent on $Z$. The same conclusions hold,
using $d=0$.

In both cases $p$ is a local minimum of $f$, by
Lemma~\ref{lem:regular-primary}(a).
\item \emph{Maximum.} Suppose $\kappa^c(\alpha,\infty)>0$, with $Z$ the set
it measures. There $C_p$ is good at $i$, hence bad at $j$. Let $F$ be silent
on $Z$. Then $F\prec_jC_p$ and $\Phi(F)=f(p)$.

Suppose instead $\kappa^d(\beta,\alpha)>0$, with $Z$ the set it measures.
There $C_p$ is silent. Let $F$ give the $j$-good advice on $Z$, which is bad
at $i$. Again $F\prec_jC_p$ and $\Phi(F)=f(p)$.

In both cases $p$ is a local maximum of $f$ on the edge, by
Lemma~\ref{lem:regular-primary}(b).
\end{itemize}
The local-minimum values of a real function on a second-countable space are
countable (CC, proof of Lemma~10), and so are its local-maximum values.

\emph{Step 3: free points.} Call $p$ \emph{free} if neither mechanism applies
at $p$. The sets $\{t=\alpha\}$ and $\{t=\beta\}$ are null. So at a free
point
\begin{equation}
 \kappa^c(\beta,\infty)=0,\qquad \kappa^d(0,\alpha)=0.
 \label{eq:free}
\end{equation}
Put
\[
 U\coloneqq\inf\{\tau\in[0,\infty]:\kappa^c(\tau,\infty)=0\},\qquad
 L\coloneqq\sup\{\tau\in[0,\infty]:\kappa^d(0,\tau)=0\},
\]
with $\kappa^c(\infty,\infty)=\kappa^d(0,0)=0$, so $U,L\in[0,\infty]$.
By continuity of $\mu$ along monotone unions,
\begin{equation}
 \kappa^c(U,\infty)=0,\qquad \kappa^d(0,L)=0.
 \label{eq:UL}
\end{equation}
These numbers do not depend on $p$. At a free point, \eqref{eq:free} gives
$U\le\beta\le\alpha\le L$. By \eqref{eq:sandwich}, every free value of $f$
lies in $[m_0+G(U),m_0+G(L)]$.

If there is no free point, every value in $[m_0,m_0+G(\infty)]$ is a local
extremum value, so this interval is countable. Hence $G(\infty)=0$, and $\Phi$
is blind on $D_{ij}$.

Otherwise, every value in
\[
 [m_0,m_0+G(U))\ \cup\ (m_0+G(L),m_0+G(\infty)]
\]
is taken only at points that are not free, so these intervals are countable.
Hence $G(U)=0$ and $G(L)=G(\infty)$, so
\begin{equation}
 c=d=0\ \text{a.e.\ on }D_{ij}\cap\{t<U\}\ \text{and on }D_{ij}\cap\{t>L\}.
 \label{eq:zero}
\end{equation}

\emph{Step 4: cases.} Recall that $t$ has a positive density on $(0,\infty)$
under $\mu|_{D_{ij}}$.
\begin{itemize}[leftmargin=*,itemsep=2pt]
\item If $U>0$ and $L>0$, then $D_{ij}\cap\{t<\min(U,L)\}$ is non-null. There
$d=0$ a.e.\ by \eqref{eq:zero}, and $d>0$ a.e.\ by \eqref{eq:UL}. This is
impossible.
\item If $L=0$, then \eqref{eq:zero} gives $c=d=0$ a.e.\ on $D_{ij}$: blind.
\item If $U=0$, then \eqref{eq:UL} gives $c>0$ a.e.\ on $D_{ij}$. By
\eqref{eq:zero}, $D_{ij}\cap\{t>L\}$ is null, so $L=\infty$. Then
\eqref{eq:UL} gives $d>0$ a.e.\ on $D_{ij}$: strict.
\end{itemize}
The two alternatives exclude each other because $D_{ij}$ is non-null.

\emph{Step 5: $n\ge3$.} Every mixed problem lies in some $D_{ij}$. For
distinct $j,k\ne i$, $D_{ij}\cap D_{ik}$ is non-null, and $c>0$ a.e.\ and
$c=0$ a.e.\ cannot both hold on it. So the alternatives on $D_{ij}$ and $D_{ik}$
agree.
\end{proof}

\begin{remark}
The theorem matches the known examples.
\begin{itemize}[leftmargin=*,itemsep=1pt]
\item Strictly ordered additive losses are strict.
\item CC's Proposition~3 ranks first by errors on unanimous problems. That
primary is blind on $H'$, and (S) holds.
\item BA's lexicographic primary has $c=0<d$ on $H'$. That is neither
alternative, and BA's Proposition~6 shows that (S) fails.
\end{itemize}
\end{remark}

At world $i$, call a ranking \emph{lexicographic over additive criteria} on
regular models if there are finitely or countably many functionals
$\Phi^1,\Phi^2,\ldots$ of the additive form above, each with weakly ordered
penalties and an integrable envelope, such that for regular $\Dscr,\Escr$:
\begin{itemize}[leftmargin=*,itemsep=1pt]
\item if $(\Phi^k(\Dscr))_k$ is lexicographically smaller than
$(\Phi^k(\Escr))_k$, then $\Dscr\prec_i\Escr$;
\item if $\Dscr\preceq_i\Escr$, then $(\Phi^k(\Dscr))_k$ is lexicographically
at most $(\Phi^k(\Escr))_k$.
\end{itemize}
Ties in the whole sequence may be broken, or left incomparable, in any way.
Call a level \emph{blind on $H'$} if its $c$ and $d$ vanish a.e.\ on $H'$.

\begin{theorem}\label{thm:lex}
Let $n\ge3$, and assume (G), (LA) and (S). Suppose that at every world the
ranking is lexicographic over additive criteria on regular models, with at
least one level not blind on $H'$. Then every $\Dscr_M$ whose credal set
contains three simplex vertices is dominated by a coherent model. In
particular the vacuous model is dominated.
\end{theorem}
\begin{proof}
Fix a world $i$, and let $\Phi^k$ be the first level not blind on $H'$.
Earlier levels are constant on regular models, since regular models are
correct on unanimous problems. So $\Phi^k$ is a regular primary: if
$\Phi^k(\Dscr)<\Phi^k(\Escr)$, the sequences are lexicographically ordered.
By Theorem~\ref{thm:dichotomy}, $\Phi^k$ is strict on $H'$.

Define $\ell_i$ on $H'$ by $\varphi^k(\cdot,r)-\varphi^k(\cdot,0)$ for
$r=0,1,2$ and by $\varphi^k(\cdot,2)-\varphi^k(\cdot,0)+1$ for contradiction.
On unanimous problems use $0,1,2,3$. This is an additive guidance loss, with
envelope $2h+3$. For regular $\Dscr$,
$L_i(\Dscr)=\Phi^k(\Dscr)-\int_H\varphi^k(\cdot,0)\,d\mu$.

If $\Dscr\preceq_i\Escr$ with both regular, the earlier levels tie, so
$\Phi^k(\Dscr)\le\Phi^k(\Escr)$ and $L_i(\Dscr)\le L_i(\Escr)$. This holds at
every world, so CC's Theorem~1 applies.
\end{proof}

So a lexicographic counterexample must be blind on mixed problems at every
level, at some world. All of its sensitivity to mixed problems then lies in
the final tie-breaking, which must still carry strict (G).

\section{One non-additive world}

The next theorem lets one world be arbitrary. It shows that the other two
worlds alone already force the vacuous trade-off on their common face.

Let $n=3$. A problem is \emph{$(2,3)$-mixed} if $g_2g_3<0$. Orient it so that
$g_2>0>g_3$; then $g_1$ may have either sign. Put
\[
 R_1\coloneqq\{g_1>0,\ g_2>0>g_3\},\qquad R_2\coloneqq\{g_1<0,\ g_2>0>g_3\}.
\]
For an additive guidance loss $L_k$, write $c_k$ and $b_k$ for its silence
and bad penalties on these problems. Accepting $g$ costs $(0,b_3)$ at worlds
$(2,3)$, silence costs $(c_2,c_3)$, and accepting $-g$ costs $(b_2,0)$. Put
\[
 \alpha\coloneqq\frac{c_2}{b_3-c_3},\qquad
 \beta\coloneqq\frac{b_2-c_2}{c_3},\qquad
 \gamma\coloneqq\frac{b_2}{b_3},\qquad
 t\coloneqq\frac{g_2}{-g_3}.
\]

\begin{theorem}\label{thm:one-world}
Let $n=3$, assume (G) and (S), and suppose the vacuous model is undominated.
Suppose worlds $2$ and $3$ have additive guidance losses $L_2,L_3$ with
\[
 L_k(C)\le L_k(X)\qquad\text{for all }(C,X)\in\Kscr_k,\ k=2,3.
\]
World $1$ is unrestricted apart from (G). Then some $\kappa^*\in(0,\infty)$
satisfies, a.e.\ on $(2,3)$-mixed problems,
\begin{equation}
 \frac{c_2}{b_2}+\frac{c_3}{b_3}\le1,\qquad
 c_2\le\kappa^*(b_3-c_3),\qquad
 \kappa^*c_3\le b_2-c_2.
 \label{eq:forced-23}
\end{equation}
Equivalently, silence is optimal on every $(2,3)$-mixed problem for the
weights $(1,\kappa^*)$ on worlds $(2,3)$.
\end{theorem}
\begin{proof}
\emph{Step 1: coherent Bayes models on the face.} Fix $\kappa\in(0,\infty)$.
Let $B$ choose, at each problem, an advice minimizing $\ell_2+\kappa\ell_3$,
with a fixed tie-break. Correct advice on unanimous problems costs $0$ and
every other advice costs more, and contradiction costs more than silence. So
$B$ is regular.

On the edge $[e_2,e_3]$, the closed sets $A_2(B)$ and $A_3(B)$ cover the
edge, by competitiveness of $C_r$ on $\supp r\subseteq\{2,3\}$. They contain
$e_2$ and $e_3$ respectively. The edge is connected, so some $r$ lies in both.
Then $(C_r,B)\in\Kscr_2\cap\Kscr_3$, so $L_k(C_r)\le L_k(B)$ for $k=2,3$, and
$E_\kappa\coloneqq C_r$ is a coherent Bayes model for the weights
$(1,\kappa)$.

Where $g_2,g_3>0$, accepting $g$ is the unique optimum for these weights. So,
by the argument of v3's Lemma~11(c), the credal set of $E_\kappa$ lies on the
edge $[e_2,e_3]$. Its advice depends only on $(g_2,g_3)$.

\emph{Step 2: the Bayes thresholds are functions of $t$.} Put
$A\coloneqq\min(\alpha,\gamma)$ and $B'\coloneqq\max(\beta,\gamma)$. At the
weights $(1,\kappa)$, accepting $g$ is strictly optimal if and only if
$\kappa<A$, and accepting $-g$ is strictly optimal if and only if
$\kappa>B'$. So, up to null sets on $R_1\cup R_2$,
\[
 \{A>\kappa\}\subseteq E_\kappa\subseteq\{A\ge\kappa\}.
\]
Since $E_\kappa$ depends only on $(g_2,g_3)$, its acceptance set on
$R_1\cup R_2$ is an upper interval in $t$, the same on $R_1$ and $R_2$. For
rational $\kappa$,
$\{A>\kappa\}=\bigcup_{\kappa'\in\mathbb Q,\,\kappa'>\kappa}E_{\kappa'}$ a.e.
The argument in the first part of v3's Lemma~12 then gives a nondecreasing
$F:(0,\infty)\to[0,\infty]$ with $A=F(t)$ a.e.\ on $R_1\cup R_2$. Rejection sets
are lower intervals in $t$, and in the same way $B'=G(t)$ a.e.\ for a
nondecreasing $G$.

\emph{Step 3: a consequence of competitiveness.} By (KKMb), the vacuous model
$V$ is competitive on some nonempty support, and therefore on $\{1,2,3\}$.
Take measurable $Q_1\subseteq R_1$ and $Q_2\subseteq-R_2$; both lie in $H$.
Let $X$ agree with $V$ except that it accepts $g$ on $Q_1\cup Q_2$. Then $X$ is
regular. At world $1$ it replaces silence by good advice on $Q_1\cup Q_2$, so
$X\prec_1V$ by strict (G) when $Q_1\cup Q_2$ is non-null. Competitiveness gives
$V\preceq_2X$ or $V\preceq_3X$. Since $V$ is coherent and undominated,
\[
 \max\bigl(L_2(X)-L_2(V),\ L_3(X)-L_3(V)\bigr)\ge0.
\]
On $Q_2$, accepting $g$ means accepting $-g'$ for the oriented $g'=-g\in R_2$.
Hence
\[
 L_2(X)-L_2(V)=-\int_{Q_1}c_2+\int_{Q_2}(b_2-c_2),\qquad
 L_3(X)-L_3(V)=\int_{Q_1}(b_3-c_3)-\int_{Q_2}c_3,
\]
where the integrands on $Q_2$ are evaluated at $-g$.

\emph{Step 4: separation of $\alpha$ and $\beta$.} We claim
\[
 \operatorname*{ess\,sup}_{R_1}\alpha\ \le\ \operatorname*{ess\,inf}_{R_2}\beta.
\]
If not, choose $0<\kappa_2<\kappa_1$ strictly between the two sides, with
$\kappa_2$ above the infimum and $\kappa_1$ below the supremum. Take non-null
$P_1\subseteq R_1\cap\{\alpha>\kappa_1\}$ and
$P_2\subseteq-(R_2\cap\{\beta<\kappa_2\})$. The measures $(b_3-c_3)\,d\mu$ and
$c_3\,d\mu$ are atomless with positive densities. So we can choose a non-null
$Q_1\subseteq P_1$ with
\[
 x\coloneqq\int_{Q_1}(b_3-c_3)\le\frac{\kappa_2}{\kappa_1}\int_{P_2}c_3,
\]
then $y\in(x,\kappa_1x/\kappa_2]$, and then $Q_2\subseteq P_2$ with
$\int_{Q_2}c_3=y$. Now $L_3(X)-L_3(V)=x-y<0$. Also
$\int_{Q_1}c_2>\kappa_1x$ because $c_2=\alpha(b_3-c_3)$, and
$\int_{Q_2}(b_2-c_2)<\kappa_2y\le\kappa_1x$ because $b_2-c_2=\beta c_3$. So
$L_2(X)-L_2(V)<0$. This contradicts Step 3.

Since $\alpha>0$ and $\beta<\infty$ a.e., we can choose $\kappa^*\in(0,\infty)$
between the two sides of the claim.

\emph{Step 5: $F\le\kappa^*\le G$ everywhere.} On $R_1$, $A\le\alpha\le\kappa^*$
a.e., so $F(t)\le\kappa^*$ for a.e.\ $t$, because $t$ has a positive density
on $R_1$. The set $\{F>\kappa^*\}$ is an upper interval of measure zero, so it
is empty. In the same way, $B'\ge\beta\ge\kappa^*$ a.e.\ on $R_2$ gives
$G\ge\kappa^*$ everywhere. This step uses only the definitions of $A$ and $B'$.

\emph{Step 6: conclusion.} The following are equivalent, since each reduces to
$b_2c_3+c_2b_3\le b_2b_3$:
\[
 \frac{c_2}{b_2}+\frac{c_3}{b_3}\le1,\qquad \alpha\le\gamma,\qquad
 \gamma\le\beta,\qquad \alpha\le\beta.
\]
Where they fail, $\alpha>\gamma>\beta$, so $A=B'=\gamma$.
\begin{itemize}[leftmargin=*,itemsep=1pt]
\item If they fail on a non-null $P\subseteq R_1$, then on $P$,
$B'=\gamma<\alpha\le\kappa^*$, contradicting $G\ge\kappa^*$.
\item If they fail on a non-null $P\subseteq R_2$, then on $P$,
$A=\gamma>\beta\ge\kappa^*$, contradicting $F\le\kappa^*$.
\end{itemize}
So the first inequality in \eqref{eq:forced-23} holds a.e., and $A=\alpha$,
$B'=\beta$ a.e. Hence $\alpha=F(t)\le\kappa^*\le G(t)=\beta$ a.e.\ on
$R_1\cup R_2$. These are the other two inequalities.
\end{proof}

Theorem~\ref{thm:one-world} recovers v3's forced edge bounds
(\emph{forced-edge-bounds} in v3, for the pair $\{2,3\}$ with
$p^0_3/p^0_2=\kappa^*$), without any numerical structure at world $1$. Strict
(G) is used only at world~$1$.

\begin{corollary}\label{cor:reduction}
Under the hypotheses of Theorem~\ref{thm:one-world}, let $c_k,b_k$ now denote
the silence and bad penalties of $L_k$ on all problems, and define, for
$k=2,3$,
\[
 \psi_k\coloneqq c_k\ \text{where }g_k>0,\qquad
 \psi_k\coloneqq-(b_k-c_k)\ \text{where }g_k<0.
\]
Let $\psi_1$ be any finite measurable function with
$\operatorname{sgn}\psi_1=\operatorname{sgn}g_1$ a.e. Then for some $p\in\Delta$
with $p_1>0$, no coherent $\Dscr_K$ with $K\subseteq\Delta_{\supp p}$ satisfies
\[
 \{p\cdot\psi>0\}\subseteq\Dscr_K\subseteq\{p\cdot\psi\ge0\}
 \quad\text{a.e.}
\]
\end{corollary}
\begin{proof}
Suppose such a $\Dscr_K$ exists for every $p$ with $p_1>0$. We show that one
also exists on the edge $[e_2,e_3]$. At $e_2$ and $e_3$ take $D_2$ and $D_3$.
For $p$ inside the edge, with $\kappa=p_3/p_2$, take $E_\kappa$ from Step~1 of
the proof above. It accepts only where acceptance is optimal, so
$E_\kappa\subseteq\{p\cdot\psi\ge0\}$. Conversely, where $p\cdot\psi>0$,
acceptance strictly beats silence. On $(2,3)$-mixed problems the first
inequality of \eqref{eq:forced-23} gives $\alpha\le\gamma\le\beta$. So
acceptance of the oriented $g$ beats silence only if $\kappa<\alpha\le\gamma$,
and rejection beats silence only if $\kappa>\beta\ge\gamma$. In both cases it
also strictly beats the third option. On $(2,3)$-unanimous problems this is
immediate. So $E_\kappa$ accepts there.

Now v3's Lemma~12 applies with $I=\{1,2,3\}$. It makes
$\psi_2/(-\psi_3)=\alpha$ unbounded on $\{g_2>0>g_3\}$. This contradicts
$\alpha\le\kappa^*$.
\end{proof}

When world $1$ also has an additive loss, v3 shows that the coherent Bayes
models at priors charging world $1$ are such completions, with $\psi_1$ the
acceptance-versus-silence contrast of that loss (v3, proof of
Proposition~13). Corollary~\ref{cor:reduction} says that, with world $1$
arbitrary, this fails at some such prior for every choice of $\psi_1$, and for
every coherent model there, selected or not.

\section{What remains}

None of the results above uses (IB). Combining them, suppose that for some
$n\ge3$ a ranking satisfies (G), (LA) and (S), and leaves the vacuous model
undominated. Then:
\begin{enumerate}[label=(\roman*),leftmargin=*,itemsep=2pt]
\item at some world $i$, the one-sided comparisons $\Kscr_i$ contain a
leverage sequence (Proposition~\ref{prop:one-sided},
Theorem~\ref{thm:leverage} and CBOI's Theorem~2);
\item at that world, the ranking on regular models is not lexicographic over
additive criteria with a level that is not blind on $H'$. Such a ranking has
an ordinary additive weak extension (proof of Theorem~\ref{thm:lex}), and so
no leverage sequence;
\item if $n=3$ and worlds $2,3$ have ordinary one-sided additive extensions,
then world $1$ is the leverage world, and for every world-$1$ contrast
$\psi_1$ some selection $C_p$ with $p_1>0$ is not a contrast completion
(Corollary~\ref{cor:reduction}).
\end{enumerate}
Item (iii) is proved only for ordinary additive extensions at the two other
worlds. We have not checked whether it extends to finitely additive ones.

The conjecture of CC therefore reduces to the following question. Since
nothing above uses (IB), an answer must use it.

\paragraph{Question.} Do (G), (IB) and (S) exclude leverage sequences in
$\Kscr_i$ at every world?

A yes proves the conjecture. A counterexample to the conjecture must give a
no at some world. The difficulty is this. (IB) transfers a comparison to
models that differ from the original pair only where the pair agree. The new
models are usually no longer coherent. So (IB) acts on comparisons between
regular models, while the extension is needed only for comparisons from
coherent undominated models. A proof would need (S) to bring the transferred
comparisons back to selected models.

\begin{thebibliography}{9}
\bibitem{Rokhlin} D.~B. Rokhlin, The Kreps--Yan theorem for $L^\infty$,
\emph{International Journal of Mathematics and Mathematical Sciences}
2005(17), 2749--2756. doi:10.1155/IJMMS.2005.2749.
\end{thebibliography}
\end{document}
